% Section source for the main manuscript; uses standard AMS math commands.
\section{The exact arithmetic foundation and the Lean endpoint}
\label{sec:foundation-audit}

\subsection{What the repository supplies}
We inspected ProveIt at the following commit:
\begin{center}\small\ttfamily
b5a3ad037723c44d69a3d5b28a950c554327a989.
\end{center}
The file \texttt{ConciseIPRA/PR.lean} defines a deeply embedded first-order
theory: its function symbols are syntax trees built from zero, successor,
projections, composition, and primitive recursion. Its nonlogical axioms
are the defining equations, separation and injectivity of successor, and
induction for quantifier-free formulas. The deduction rules in
\texttt{Deduction.lean} are intuitionistic natural deduction with equality.
The notation
\[
       \Gamma\vdash[\mathsf{PRAAx}]\varphi
\]
denotes a derivation in that object theory; it does not merely assert that
the standard natural numbers satisfy~$\varphi$.

The native scalar calculus uses iterator terms and formulas consisting
only of equations and two binding constructions. Its ordinary readings are
\[
 Q_yB(y)\;\equiv\;B(0)\to\forall y\,B(y),\qquad
 D_{x,y}B(x,y)\;\equiv\;
     \exists x\bigl(B(x,0)\to\forall y\,B(x,y)\bigr).
\]
The inspected theorem \texttt{pra\_iff\_native}, in
\texttt{NativeBiInterpretation.lean}, says exactly that
\[
 (\Gamma\vdash[\mathsf{PRAAx}]\varphi)
 \quad\Longleftrightarrow\quad
 \mathsf{NPrf}\bigl(\Gamma.\mathsf{map}\,\mathsf{NFm.ofPRA}\bigr)
                   (\mathsf{NFm.ofPRA}\,\varphi).
\]
Both interpretations preserve the whole numerical domain and the
repository also proves both formula round trips. Consequently the natural
formalization route is to prove the finite arithmetic estimates in the
ordinary iPRA presentation, then apply this established bridge. Rebuilding
the smallest native proof checker is not a prerequisite for this work.

\subsection{Finite statements rather than hidden higher-order quantifiers}
Every rational is coded by a signed numerator and a positive denominator;
complex rationals are pairs of rationals. Finite lists, polynomial
coefficients, meshes, contour samples, and finite certificates are numerical
codes. Their comparisons are decidable by integer arithmetic. Malformed
codes receive a fixed default output, so each implementation is a total
primitive recursive function. Formula quantifiers range over such natural
codes, not over arbitrary real numbers, analytic functions, or primitive
recursive programs.

For a fixed Boolean checker $C$ with $r$ numerical inputs, its internal
correctness statement is
\[
         \forall x_1\cdots\forall x_r\;C(x_1,\ldots,x_r)=1.
\]
The quantifier-free matrix is an equation in the language of iPRA. A
conditional checker returns $1$ when its finite input hypotheses are false
and otherwise checks the asserted rational conclusion. This convention
does not dispense with hypotheses: the report must specify exactly which
hypotheses the finite checker recognizes.

Statements for a primitive recursive Cauchy name are schemas in a
\emph{fixed definition} of that name. Its regularity hypotheses may remain
universally quantified assumptions. No step evaluates all primitive
recursive program codes by a single total primitive recursive function.
That distinction is essential: a universal evaluator of that kind is not
available.

At each use of induction the requested precision and all size parameters
are fixed first. The invariant is an equation or inequality about rational
partial sums, products, mesh-cell errors, or a bounded Boolean fold. A
bounded conjunction $\bigwedge_{j<n}P(j)$ is represented by the product of
the characteristic values of the fixed predicate~$P$. An induction over
such a fold is quantifier-free induction. This does not authorize induction
over an assertion such as ``the analytic equality holds at every
precision'', or over the existence of an unspecified zero. Those larger
claims are assembled only after the finite inductions have finished.

\subsection{What a contour certificate certifies}
A finite contour certificate records a rational polygon in the domain,
rational samples of the fixed analytic approximation, a numerical sample
precision, a derivative or continuity bound, a positive rational margin,
and an integer winding count. Its checker recomputes the samples and tests
the error and separation inequalities. Sample values are not trusted
merely because the certificate contains them.

For example, if the approximation error is at most $e$, each domain edge
has length at most $h$, and the derivative bound is~$L$, then an inequality
of the form
\[
            |w_j|_\infty>4e+4Lh
\]
at every sample gives a zero-free homotopy from the actual boundary image
to its polygonal approximation, after including the corresponding
successive-sample bounds. All tests can use rational norms; no square-root
oracle is required. The winding of the rational polygon is computed by
oriented ray crossings. A transverse rational ray can be selected by a
finite search which excludes the finitely many vertex directions.

There are two separate formal obligations. First, a valid contour has the
claimed count and obeys the distribution estimate. Second, the specified
primitive recursive construction produces valid contours at the stated
heights. The first obligation alone is insufficient for unconditional
unboundedness of the zero count. The minimum-modulus and regular-height
construction in the counting argument exists precisely to supply the
second obligation. Nor is a certified count the same as an algorithm that
sorts exact zeros by ordinate or decides equal ordinates.

\subsection{The inspected analytic drafts and their status}
The repository's archived functional-equation draft supplies a quantitative
human proof of the theta transformation, reciprocal gamma products,
reflection and duplication, and the Mellin identification with the zeta
Dirichlet series. Its finite xi approximation has exact algebraic symmetry;
its identification with the analytic zeta function requires the substantive
Mellin proof. The archive explicitly says that its Lean interfaces are
uncompiled and incomplete. We reuse its written estimates at that stated
level, not as a kernel-checked theorem.

The Cauchy--residue draft gives finite triangulation cancellation,
quantitative Laurent expansions and rational contour winding calculations.
Its core singularity locations are supplied rational centres. A generic
argument principle for unknown zeta zeros does not follow merely by citing
that theorem: the local factorization, counting and approximation work
must also be supplied. Similarly, the Stirling draft proves the real
factorial formula with explicit remainders; it is not a proof of complex
log-gamma asymptotics. The complex estimate used in zero counting needs its
own proof. The FTA draft provides a written quantitative construction of
primitive recursive root names, but also explicitly lacks a completed
object-language Lean derivation.

\subsection{A concrete completion order}
The accompanying \texttt{ObjectTargets.lean} is an uncompiled interface.
It adds no analytic axiom and supplies no instance of its final obligation
structure. The implementation sequence is as follows.
\begin{enumerate}
\item Compile integer, rational, list, bounded-search and polynomial
      operations to the repository's \texttt{PR} syntax. Prove defining
      equations and the required arithmetic lemmas inside
      \texttt{Prf PRAAx}; semantic Lean calculations alone do not suffice.
\item Compile the fixed elementary and zeta/xi approximation programs.
      Formalize their rational error budgets, input sensitivities and
      derivative bounds with the displayed quantifier-free invariants.
\item Prove the finite functional-equation, Euler-product, and gamma
      identities. Keep numerical division guarded by an explicit positive
      lower bound. Prove the positive nonvanishing estimates.
\item Implement rational contour certification and winding. Prove
      finite homotopy invariance and the counting statements, including
      the multiplicity convention.
\item Formalize the quantitative Jensen, minimum-modulus, complex
      Stirling and argument-variation estimates. Compile the bounded
      regular-height generator and prove that its output passes the
      certificate checker.
\item Construct the internal derivations of the rational zero-counting
      and reciprocal-tail sentences. Derive unboundedness and the stated
      distribution corollaries from these internal finite estimates.
\item Apply \texttt{pra\_iff\_native} to each completed derivation.
      Kernel-check all resulting proof terms at the pinned toolchain,
      and audit dependencies for \texttt{sorryAx} and additional axioms.
\end{enumerate}
The inspected repository requests Lean and mathlib \texttt{v4.32.0}.
Those versions identify a proposed build environment; this package does not
claim that its new analytic development has been compiled there. A proof
about \texttt{Complex} in ambient Lean can be used as an independent
semantic comparison, but does not establish the restriction to iPRA.
